% `tex.web` §1064 `off_save`, and §1065's table of what it inserts.
%
%     procedure off_save;
%     var p:pointer; {inserted token}
%     begin if cur_group=bottom_level then
%       @<Drop current token and complain that it was forbidden@>
%     else  begin back_input; p:=get_avail; link(temp_head):=p;
%       print_err("Missing ");
%       @<Prepare to insert a token that matches |cur_group|@>;
%       print(" inserted"); ins_list(link(temp_head));
%       help5(…); error;
%       end;
%     end;
%
%     case cur_group of
%     semi_simple_group: … print_esc("endgroup");
%     math_shift_group:  … print_char("$");
%     math_left_group:   … print_esc("right.");
%     othercases         … print_char("}")
%
% A command arrived that cannot be obeyed with the current group open, so the
% token that *would* close the group is inserted in front of it. The inserted
% token is then executed, which is what makes the choice matter: the wrong one
% closes the wrong thing and everything after diverges.
%
% We had the `}` arm only, written inline at its one caller (§1095's
% `head_for_vmode`). §1066's other half was missing too — with *no* group open
% there is nothing to insert, so the token is dropped and named:
%
%     begin print_err("Extra "); print_cmd_chr(cur_cmd,cur_chr);
%     help1("Things are pretty mixed up, but I think the worst is over.");
%     error;
%     end
%
% and the two callers wired up here are §1065's `\endgroup` — legal only for a
% `\begingroup`'s own group — and §1145's `\halign`, which in a display is
% legal only when the group is the display's.
%
% Not covered: the `\right.` arm. It is implemented and reachable, but every
% way of reaching it also runs through §1193's `Extra \right`, which does not
% exist yet, so a case for it would be testing that instead.
\catcode`\{=1 \catcode`\}=2 \catcode`\$=3 \catcode`\#=6 \catcode`\&=4
\errorcontextlines=5
\hsize=100pt \parindent=0pt
\font\rm=cmr10 \textfont0=\rm \scriptfont0=\rm \scriptscriptfont0=\rm \rm
\font\mi=cmmi10 \textfont1=\mi \scriptfont1=\mi \scriptscriptfont1=\mi
\font\sy=cmsy10 \textfont2=\sy \scriptfont2=\sy \scriptscriptfont2=\sy
\font\ex=cmex10 \textfont3=\ex \scriptfont3=\ex \scriptscriptfont3=\ex
\showthe\escapechar

% The `}` arm, twice over: the inserted brace closes the `{`, the `\endgroup`
% comes back and closes the \hbox the same way, and only then is there no
% group left and §1066 drops it.
\setbox0\hbox{A{B\endgroup C}D}
\message{PROBE:brace=[\the\wd0]}

% The `$` arm. The cascade is the point: the inserted `$` closes the formula,
% the `\endgroup` comes back and meets the \hbox, and §1047 picks up the
% unclosed formula the trailing `$D}` leaves behind.
\setbox1\hbox{A$B\endgroup C$D}
\message{PROBE:math=[\the\wd1]}

% §1145's caller: an alignment in a display, inside a group that is not the
% display's. This is trip.tex line 250.
\setbox2\vbox{\noindent A$$\begingroup\halign{#\cr B\cr}$$C\par}
\message{PROBE:halign=[\the\ht2]}

% ...and the same alignment with no group in the way, which is legal.
\setbox3\vbox{\noindent A$$\halign{#\cr B\cr}$$C\par}
\message{PROBE:plain=[\the\ht3]}

% §1066 on its own: nothing open at all.
\endgroup
\message{PROBE:bottom}

% A `\begingroup` closed by its own `\endgroup` says nothing.
\begingroup\endgroup
\message{PROBE:ok}
\showthe\escapechar
\end
